% !TEX root = ../main.tex
\questiontitle{Question 1: IoT Sensing Network over 1 km $\times$ 1 km}

\begin{question}{1}
Consider an IoT sensing network to be deployed over a 1 km $\times$ 1 km (1000 m $\times$ 1000 m)
square area. A single Internet-connected Access Point (AP) is located at the top-right corner of the sensing area.

Assumptions:
\begin{enumerate}[label=\arabic*.]
\item Each IoT sensor node has a maximum wireless communication range of 100 m. Two IoT nodes can directly communicate if their Euclidean separation is $\leq$ 100 m.
\item The AP (Access Point) has a wireless communication range of 200 m.
\item Each IoT node generates one sensor-data packet every minute.
\item IoT nodes that cannot directly communicate with the AP (Access Point) must forward their
data through other IoT nodes using multi-hop communication.
\item Assume ideal wireless conditions: no packet loss, interference, retransmissions, or physical obstacles.
\item The objective is to achieve complete network connectivity to the AP (Access Point) while using as few IoT nodes as reasonably possible.
\item You are not given a predefined grid spacing or node coordinates. You must determine and justify an appropriate
placement.
\end{enumerate}
\end{question}

% ------------------------------------------------------------------ (a)
\begin{question}{1 (a)}
Design a suitable placement of IoT nodes over the 1 km $\times$ 1 km (1000 m $\times$ 1000 m)
area. Determine an appropriate node spacing and the total number of nodes required. Draw the proposed placement,
clearly showing the AP (Access Point) and the IoT nodes.
\end{question}

\answerhead{(a) Placement, spacing and number of nodes}
\begin{itemize}
\item Square grid with node spacing \textbf{d = 100 m}.
\item First node N1 at \textbf{(50, 50)}.
\item This gives 10 $\times$ 10 = 100 grid positions. The last top-right position \textbf{(950, 950)} is the
\textbf{AP (Access Point)}.
\item Total IoT nodes required = \textbf{99} (N1 to N99, numbered row by row from the
bottom-left).
\end{itemize}

\simfigure[width=0.86\linewidth]{../media/figures/q1_layout.pdf}{Proposed placement: 99 IoT
nodes at 100 m spacing and the AP (Access Point) at (950, 950). Grey nodes are within the 200 m range of the AP
(Access Point).}

% ------------------------------------------------------------------ (b)
\begin{question}{1 (b)}
Justify your selected node spacing and placement. Demonstrate that every IoT node can ultimately
reach the AP (Access Point), either directly or through multi-hop communication, while maintaining the $\leq$ 100 m
communication-distance constraint between directly communicating IoT nodes.
\end{question}

\answerhead{(b) Justification and connectivity}
\begin{itemize}
\item \textbf{Spacing:} neighbouring nodes (left, right, up, down) are exactly 100 m apart, which equals the node
range, so every neighbour link satisfies the $\leq$ 100 m constraint.
\item \textbf{Fewest nodes:} any spacing larger than 100 m breaks the links; any smaller spacing adds nodes that are
not needed. 100 m is therefore the largest valid spacing and needs the fewest nodes.
\item \textbf{Coverage:} the 50 m offset from the edges places each node at the centre of its own 100 m $\times$ 100
m cell, so no point of the area is more than 70.7 m from a node.
\item \textbf{Every node reaches the AP (Access Point):} from any node, the neighbour 100 m to its right or 100 m
above it is one hop closer to the AP (Access Point). Repeating this step always reaches one of the five nodes inside
the 200 m range of the AP (Access Point), which then sends directly to the AP (Access Point). The farthest node, N1,
needs 17 transmissions; all 99 nodes are connected.
\end{itemize}

% ------------------------------------------------------------------ (c)
\begin{question}{1 (c)}
Identify the IoT nodes that can communicate directly with the AP (Access Point). Select a node
that cannot directly reach the AP (Access Point) and determine a shortest multi-hop path from that node to the AP
(Access Point). State the number of IoT-to-IoT hops and total wireless
transmissions, including the final AP (Access Point) transmission.
\end{question}

\answerhead{(c) Direct nodes and shortest path}
\textbf{Nodes in direct range of the AP (Access Point) at (950, 950), distance $\leq$ 200 m:}
\begin{itemize}
\item N99 at (850, 950): 100 m
\item N90 at (950, 850): 100 m
\item N89 at (850, 850): 141.4 m
\item N98 at (750, 950): 200 m
\item N80 at (950, 750): 200 m
\end{itemize}

\textbf{Selected node:} N1 at (50, 50).

\textbf{Shortest path:} N1 $\rightarrow$ N11 $\rightarrow$ N21 $\rightarrow$ N31 $\rightarrow$ N41 $\rightarrow$ N51
$\rightarrow$ N61 $\rightarrow$ N71 $\rightarrow$ N81 $\rightarrow$ N82 $\rightarrow$ N83 $\rightarrow$ N84
$\rightarrow$ N85 $\rightarrow$ N86 $\rightarrow$ N87 $\rightarrow$ N88 $\rightarrow$ N89 $\rightarrow$ AP (Access
Point)

\begin{itemize}
\item IoT-to-IoT hops = \textbf{16}
\item Total wireless transmissions = 16 + 1 (final hop N89 $\rightarrow$ AP (Access Point)) = \textbf{17}
\end{itemize}

\simfigure[width=0.86\linewidth]{../media/figures/q1_path.pdf}{Shortest multi-hop path from N1 to the AP (Access
Point): 16 IoT-to-IoT hops and 17 transmissions.}

% ------------------------------------------------------------------ (d)
\begin{question}{1 (d)}
Each IoT node generates one packet per minute. Explain how traffic accumulates toward the AP
(Access Point). Identify the nodes or region expected to experience the highest forwarding load and explain the
energy-hotspot/energy-hole problem.
\end{question}

\answerhead{(d) Traffic accumulation and the energy hole}
Each node generates one packet per minute. Nodes closer to the AP (Access Point) forward their own packets as well as packets
received from nodes farther away.

The forwarding load increases toward the AP (Access Point) because traffic from multiple branches converges near the top-right
corner. Nodes adjacent to the AP (Access Point) therefore carry a disproportionately high forwarding burden.

This creates the energy-hole problem: heavily loaded relay nodes deplete their batteries faster than nodes farther
away. When these relays fail, distant nodes may lose their routes to the AP (Access Point), reducing network coverage and
lifetime.

\simfigure[width=0.86\linewidth]{../media/figures/q1_load.pdf}{Forwarding load of every node (node number on top, packets transmitted per minute below). Arrows show the traffic converging on the AP (Access Point).}

% ------------------------------------------------------------------ (e)
\begin{question}{1 (e)}
Suggest one technique that could reduce the energy burden on heavily loaded nodes and explain how it can improve the
network lifetime.
\end{question}

\answerhead{(e) Technique: unequal clustering}
Unequal clustering can reduce the energy burden on nodes near the AP (Access Point). Smaller clusters can be formed near the AP (Access Point),
reducing the number of local packets that heavily loaded nodes must process and forward. Cluster heads can aggregate
readings before forwarding them.

This reduces redundant transmissions and balances energy consumption, potentially extending the network's
operational lifetime.
